\documentclass[11pt]{article}
\usepackage{amssymb}
\usepackage{amsthm}
\usepackage{enumitem}
\usepackage{amsmath, physics}
\usepackage{bm}
\usepackage{adjustbox}
\usepackage{mathrsfs}
\usepackage{graphicx}
\usepackage{siunitx}
\usepackage[mathscr]{euscript}

\title{\textbf{Solved selected problems of Classical Electrodynamics - Hans Ohanian}}
\author{Franco Zacco}
\date{}

\addtolength{\topmargin}{-3cm}
\addtolength{\textheight}{3cm}

\newcommand{\N}{\mathbb{N}}
\newcommand{\Z}{\mathbb{Z}}
\newcommand{\Q}{\mathbb{Q}}
\newcommand{\R}{\mathbb{R}}
\newcommand{\diam}{\text{diam}}
\newcommand{\cl}{\text{cl}}
\newcommand{\bdry}{\text{bdry}}
\newcommand{\inter}{\text{int}}
\newcommand{\uvi}{\bm{i}}
\newcommand{\uvj}{\bm{j}}
\newcommand{\uvk}{\bm{k}}
\newcommand{\hatx}{\bm{\hat{x}}}
\newcommand{\haty}{\bm{\hat{y}}}
\newcommand{\hatz}{\bm{\hat{z}}}
\newcommand{\hatrho}{\bm{\hat{\rho}}}
\newcommand{\hatphi}{\bm{\hat{\phi}}}
\newcommand{\hatr}{\bm{\hat{r}}}
\newcommand{\hatn}{\bm{\hat{n}}}
\newcommand{\hattheta}{\bm{\hat{\theta}}}
\newcommand{\sci}[1]{\times 10^{#1}}
\newcommand\numberthis{\addtocounter{equation}{1}\tag{\theequation}}

\theoremstyle{definition}
\newtheorem*{solution*}{Solution}
\renewcommand*{\proofname}{\bf{Solution}}

\begin{document}
\maketitle
\thispagestyle{empty}

\section*{Chapter 7 - Vector Calculus in Spacetime}

\subsection*{Problems}
\begin{proof}[\proofname\ 2]\leavevmode\par
\begin{itemize}
\item [(a)] Let $A(\mu,\nu)$ be an object with two indices and suppose that
$A(\mu, \nu)B^\mu$ is a vector when $B^\mu$ is a vector. \\
Let us consider the transformation law of the vector $A(\mu,\nu)B^\mu$
\begin{align*}
    A'(\mu,\nu)B'^\mu = {a_\nu}^\alpha A(\mu,\alpha)B^\mu
\end{align*}
Since $B^\mu$ is also a vector then we know that 
\begin{align*}
    B^\mu = {a_\sigma}^\mu B'^\sigma
\end{align*}
Hence
\begin{align*}
    A'(\mu,\nu)B'^\mu &= {a_\nu}^\alpha A(\mu,\alpha)B^\mu
    = {a_\sigma}^\mu {a_\nu}^\alpha A(\mu,\alpha) B'^\sigma
\end{align*}
Renaming the dummy indices on the RHS $\sigma \leftrightarrow \mu$ we get that
\begin{align*}
    A'(\mu,\nu)B'^\mu 
    &= {a_\mu}^\sigma {a_\nu}^\alpha A(\sigma,\alpha) B'^\mu
\end{align*}
Since $B'^\mu$ is a vector, it can be the vector $B'^\mu = (1, 0, 0, 0)$
and hence the equation becomes 
\begin{align*}
    A'(t,\nu)
    &= {a_t}^\sigma {a_\nu}^\alpha A(\sigma,\alpha)
\end{align*}
For the vector $B'^\mu = (0, 1, 0, 0)$ we get that
\begin{align*}
    A'(x,\nu)
    &= {a_x}^\sigma {a_\nu}^\alpha A(\sigma,\alpha)
\end{align*}
The same can be proven for the vectors $B'^\mu = (0,0,1,0)$ and 
$B'^\mu = (0,0,0,1)$.\\
With this selection of vectors we see that $A(\mu, \nu)$ transforms like a
tensor, and since this must hold for every vector, then it must be that
$A(\mu, \nu)$ is a tensor.

\item [(b)] Let $A(\mu,\nu)$ be an object with two indices and suppose that
$A(\mu, \nu)C^{\mu\nu}$ is a scalar when $C^{\mu\nu}$ is a tensor. \\
Scalars are invariant under transformation laws so
\begin{align*}
    A'(\mu, \nu)C'^{\mu\nu} = A(\mu, \nu)C^{\mu\nu}
\end{align*}
Since $C^{\mu\nu}$ is a tensor then we know that
\begin{align*}
    C^{\mu\nu} = {a_\alpha}^\mu{a_\beta}^\nu C'^{\alpha\beta}
\end{align*}
Then
\begin{align*}
    A'(\mu, \nu)C'^{\mu\nu} = {a_\alpha}^\mu{a_\beta}^\nu A(\mu, \nu) C'^{\alpha\beta}
\end{align*}
Renaming the dummy indices on the RHS $\alpha \leftrightarrow \mu$ 
and $\beta \leftrightarrow \nu$, we get that
\begin{align*}
    A'(\mu, \nu)C'^{\mu\nu} = {a_\mu}^\alpha{a_\nu}^\beta A(\alpha, \beta) C'^{\mu\nu}
\end{align*}
As in part (a), if we select
\begin{align*}
    C'^{\mu\nu} = \begin{pmatrix}
        1 & 0 & 0 & 0 \\
        0 & 0 & 0 & 0 \\
        0 & 0 & 0 & 0 \\
        0 & 0 & 0 & 0
    \end{pmatrix}
\end{align*}
Then the equation becomes
\begin{align*}
    A'(t, t) = {a_t}^\alpha{a_t}^\beta A(\alpha, \beta)
\end{align*}
The same can be shown for every component of $A(\mu,\nu)$, setting to 1 only
the component we are interested in, in the tensor $C'^{\mu\nu}$. \\
With this selection of tensors we see that $A(\mu, \nu)$ transforms like a
tensor, and since this must hold for every tensor, then it must be that
$A(\mu, \nu)$ is a tensor.
\end{itemize}
\end{proof}

\cleardoublepage
\begin{proof}[\proofname\ 3]
Let $A^\mu$ and $B_{\mu}$ be two vectors such that $A^\mu B_\mu = 0$
\begin{itemize}
\item [(a)] Suppose $A^\mu A_\mu > 0$. We note that we can normalize
$A^\mu A_\mu$ such that $A^\mu A_\mu = 1$, so we may take $A^\mu$ with three
other vectors $x^\mu, y^\mu$ and $z^\mu$ as a orthonormal basis.\\
Let us suppose $A^\mu$ is the timelike vector and the other three are spacelike
vectors. \\
Since $A^\mu B_\mu = 0$ then must be that $B^\mu$ is orthogonal to $A^\mu$
and hence it must be a linear combination of the other three vectors i.e.
\begin{align*}
    B^\mu = a x^\mu + b y^\mu + c z^\mu
\end{align*}
Where $a,b,c$ are scalars. Then
\begin{align*}
    B^\mu B_\mu
    &= -(a^2 x^1x_1 + b^2 y^1y_1 + c^2 z^1z_1)
    -(a^2 x^2x_2 + b^2 y^2y_2 + c^2 z^2z_2) \\
    &\quad -(a^2 x^3x_3 + b^2 y^3y_3 + c^2 z^3z_3) \\
    &= -(a^2 + b^2 + c^2) -(a^2 + b^2 + c^2) -(a^2 + b^2 + c^2)
\end{align*}
Where we used that $x^iy_i = x^iz_i = y^iz_i = 0$ and that 
$x^ix_i = y^iy_i = z^iz_i = 1$ for $i \in \{1,2,3\}$.
Therefore $B^\mu B_\mu < 0$.

\cleardoublepage
\item [(b)] Suppose $A^\mu A_\mu = 0$ and let us suppose $B^\mu B_{\mu} > 0$
we want to arrive at a contradiction.\\
Let $B^\mu = (1, 0, 0, 0)$ since $A^\mu B_{\mu} = 0$ we get that $A^0 = 0$, 
also since $A^\mu A_\mu = 0$ we get that
\begin{align*}
    (A^0)^2 - (A^1)^2 - (A^2)^2 - (A^3)^2 = 0 \\
    (A^1)^2 + (A^2)^2 + (A^3)^2 = 0
\end{align*}
This is a contradiction since it cannot be unless $A^\mu = (0,0,0,0)$ which is
a trivial solution and for which $B^\mu$ can be any vector. \\
Therefore must be that $B^\mu B_\mu \leq 0$.\\
If $B^\mu B_\mu = 0$ then $B^\mu$ is a lightlike vector as $A^\mu$. \\
Let us take a reference frame such that $A^{\mu} = (1,1,0,0)$.\\
Since we know that $A^\mu B_\mu = A'^\mu B'_\mu$ when $A'^\mu$
and $B'_\mu$ are in a different reference frame, then taking $B^\mu$ in the
same reference frame where $A^{\mu} = (1,1,0,0)$, must be that 
\begin{align*}
    A^\mu B_\mu &= 0 \\
    A^0B^0 - A^1B^1 &= 0 \\
    B^0 - B^1 &= 0
\end{align*}
Then $B^0 = B^1$, but also we know that $B^\mu B_\mu = 0$ so
\begin{align*}
    (B^0)^2 - (B^1)^2 - (B^2)^2 - (B^3)^2 &= 0 \\
    - (B^2)^2 - (B^3)^2 &= 0
\end{align*}
Which implies that in this frame $B^\mu$ must be of the form 
$$B^\mu = (B^0, B^0, 0, 0)$$
Therefore must be that $B^\mu$ is porportional to $A^\mu$ i.e. $B^\mu = kA^\mu$.
\end{itemize}
\end{proof}

\cleardoublepage
\begin{proof}[\proofname\ 7]\leavevmode\par
\begin{itemize}
\item [(a)] We know that the relativistic momentum is given by
\begin{align*}
    \bm{p} = \frac{m\bm{v}}{\sqrt{1 - v^2/c^2}}
\end{align*}
Since the acceleration and the velocity are always parallel then we may define
our coordinate system such that 
\begin{align*}
    p_x = \frac{mv_x}{\sqrt{1 - v_x^2/c^2}}
\end{align*}
Taking the derivative with respect to time, we get that
\begin{align*}
    \dv{p_x}{t}
    &= \frac{mv_x^2}{c^2(1 - v_x^2/c^2)^{3/2}}\dv{v_x}{t}
    + \frac{m}{\sqrt{1 - v_x^2/c^2}}\dv{v_x}{t} \\
    &= \frac{m}{(1 - v_x^2/c^2)^{3/2}}\dv{v_x}{t}
    \bigg(\frac{v_x^2}{c^2} 
    + \frac{(1 - v_x^2/c^2)^{3/2}}{\sqrt{1 - v_x^2/c^2}}\bigg) \\
    &= \frac{m}{(1 - v_x^2/c^2)^{3/2}}\dv{v_x}{t}
    \bigg(\frac{v_x^2}{c^2} + 1 - \frac{v_x^2}{c^2}\bigg) \\
    &= \frac{m}{(1 - v_x^2/c^2)^{3/2}}\dv{v_x}{t}
\end{align*}
Since when the acceleration and the velocity are parallel we can always select
the coordinate system such that they are in the same axis
(in this case the $x$ axis) then we can write that
\begin{align*}
    \dv{\bm{p}}{t} &= \frac{m}{(1 - v^2/c^2)^{3/2}}\dv{\bm{v}}{t}
\end{align*}
Therefore, comparing with the equation $d\bm{p}/dt = m(v)d\bm{v}/dt$
we see that
\begin{align*}
    m(v) = \frac{m}{(1 - v^2/c^2)^{3/2}}
\end{align*}
\item [(b)] Let us consider now a particle in uniform circular motion such that
the acceleration and the velocity are always perpendicular.\\
Given that the magnitude of the velocity is constant in a uniform circular
motion when taking the derivative with respect to time of 
\begin{align*}
    \bm{p} = \frac{m\bm{v}}{\sqrt{1 - v^2/c^2}}
\end{align*}
We ge that
\begin{align*}
    \dv{\bm{p}}{t} = \frac{m}{\sqrt{1 - v^2/c^2}}\dv{\bm{v}}{t}
\end{align*}
Therefore $m(v)$ must be
\begin{align*}
    m(v) = \frac{m}{\sqrt{1 - v^2/c^2}}
\end{align*}
\item [(c)] Finally, let us consider the case where the acceleration is
neither parallel nor perpendicular to the velocity.\\
Taking the derivative with respect to time of the relativistic momentum where
$\bm{v}$ points in an arbitrary direction, we get that
\begin{align*}
    \dv{\bm{p}}{t}
    = m\dv{t}(\frac{1}{\sqrt{1 - v^2/c^2}})\bm{v}
    + \frac{m}{\sqrt{1 - v^2/c^2}}\dv{\bm{v}}{t}
\end{align*}
We note that
\begin{align*}
    \dv{t}(\frac{1}{\sqrt{1 - (\bm{v}\cdot\bm{v})/c^2}})
    &= -\frac{1}{2}\bigg(1 - \frac{v^2}{c^2}\bigg)^{-3/2}
    \dv{t}(1 - \frac{\bm{v}\cdot\bm{v}}{c^2}) \\
    &= \frac{1}{2}\bigg(1 - \frac{v^2}{c^2}\bigg)^{-3/2}
    \frac{2}{c^2}\bigg(\bm{v}\cdot\dv{\bm{v}}{t}\bigg) \\
    &= \frac{\bm{v}\cdot\dv{\bm{v}}{t}}{c^2(1 - v^2/c^2)^{3/2}}
\end{align*}
So
\begin{align*}
    \dv{\bm{p}}{t}
    = \frac{m\bm{v}\cdot\dv{\bm{v}}{t}}{c^2(1 - v^2/c^2)^{3/2}}\bm{v}
    + \frac{m}{\sqrt{1 - v^2/c^2}}\dv{\bm{v}}{t}
\end{align*}
Therefore, $d\bm{p}/dt$ is not going to have the direction of the
velocity nor the direction of the acceleration, but a combination of both and
hence there is no scalar factor $m(v)$ we can come up with such that
$$\dv{\bm{p}}{t} = m(v)\dv{\bm{v}}{t}$$
\end{itemize}
\end{proof}

\cleardoublepage
\begin{proof}[\proofname\ 8]
Since the proton is at rest, then in momentum space the vector representing the
proton is vertical and points to the vertex of the mass shell
(or the hyperboloid), this is the point $(p^0, 0, 0, 0)$.\\
If we assume the tachyon with $v = \infty$ is hitting the proton coming in
the direction of the $p^x$-axis then the final state of the proton must be
a vector with coordinates $(p^0, p^x, 0, 0)$ where $p^x \neq 0$, but we see
that there is no point of the hyperboloid that matches this constraint.\\
Therefore, the proton cannot absorb this tachyon.
\end{proof}

\cleardoublepage
\begin{proof}[\proofname\ 9]\leavevmode\par
\begin{itemize}
\item [(a)] Let a extremely relativistic particle where $v \approx c$.
This particle has an energy of
\begin{align*}
    E &= \frac{mc^2}{\sqrt{1 - v^2/c^2}}
\end{align*}
Given that $v \approx c$ we can define a small quantity $\epsilon = 1 - v/c$
such that 
\begin{align*}
    \frac{v^2}{c^2} &= (\epsilon - 1)^2 = \epsilon^2 - 2\epsilon + 1 \approx - 2\epsilon + 1
\end{align*}
So
\begin{align*}
    1 - \frac{v^2}{c^2} &= 1 + 2\epsilon - 1 = 2\epsilon
\end{align*}
Therefore
\begin{align*}
    E &= \frac{mc^2}{\sqrt{2\epsilon}} \\
    \epsilon &= \frac{1}{2}\left(\frac{mc^2}{E}\right)^2 \\
    \frac{1}{c}(c - v)
    &= \frac{1}{2}\left(\frac{mc^2}{E}\right)^2 \\
    c - v
    &= \frac{c}{2}\left(\frac{mc^2}{E}\right)^2
\end{align*}
\item [(b)] Given a proton of $1000~GeV$ the velocity difference is
\begin{align*}
    c - v
    &= \frac{299\sci{6}~m/s}{2}\left(\frac{938.27~MeV}{1\sci{6}~MeV}\right)^2
    = 131.6~m/s
\end{align*}
\item [(c)] Given an electron of $40~GeV$ the velocity difference is
\begin{align*}
    c - v
    &= \frac{299\sci{6}~m/s}{2}\left(\frac{0.511~MeV}{4\sci{4}~MeV}\right)^2
    = 0.0243~m/s
\end{align*}
\end{itemize}
\end{proof}

\cleardoublepage
\begin{proof}[\proofname\ 10]
Let us consider a particle of mass $m_1 = 1~g$ moving in the positive $x$
direction with a velocity of $v_1 = 0.9$ and another particle of mass
$m_2 = 1~g$ moving in the negative $x$ direction with velocity $v_2 = -0.99$.
We assume $c = 1$. Then the total relativistic momentum of  the system is
\begin{align*}
    P &= \frac{m_1 v_1}{\sqrt{1 - v_1^2}} + \frac{m_2 v_2}{\sqrt{1 - v_2^2}} \\
    &= 2.06 - 7.01 \\
    &= -4.94
\end{align*}
Also, the total energy of the system is
\begin{align*}
    E &= \frac{m_1}{\sqrt{1 - v_1^2}} + \frac{m_2}{\sqrt{1 - v_2^2}} \\
    &= 2.29 + 7.08 \\
    &= 9.37
\end{align*}
So the velocity of the center of momentum from the laboratory frame is
\begin{align*}
    v_{C.M.} = \frac{P}{E} = \frac{-4.94}{9.37} = -0.52
\end{align*}
On the other hand, the velocity of the center of mass is
\begin{align*}
    v_{C.mass} = \frac{m_1v_1 + m_2v_2}{m_1 + m_2} = \frac{0.9 - 0.99}{2}
    = - 0.045
\end{align*}
Therefore we see that for relativistic particles $v_{C.M.} \neq v_{C.mass}$.


\end{proof}

\cleardoublepage
\begin{proof}[\proofname\ 20]\leavevmode\par
\begin{itemize}
\item [(a)] From the conservation of energy and momentum in the instantaneous
rest frame $x'y'z't'$ of the rocket we have that
\begin{align*}
    dv' &= \frac{v_{ex}}{M}\frac{dm}{\sqrt{1 - v_{ex}^2/c^2}}
    = - \frac{v_{ex}}{M} dM
\end{align*}
Also, from the velocity-combination formula we can approximate it as
\begin{align*}
    v + dv
    &= \frac{v + dv'}{1 + vdv'/c^2} \\
    &\approx (v + dv')\bigg(1 - \frac{vdv'}{c^2}\bigg) \\
    &\approx v + dv' - \frac{v^2dv'}{c^2} \\
    &= v + \bigg(1 - \frac{v^2}{c^2}\bigg)dv'
\end{align*}
Where we used the binomial approximation and we dropped the $dv'^2$ term.
Then
\begin{align*}
    dv &= \bigg(1 - \frac{v^2}{c^2}\bigg)dv'
\end{align*}
Replacing $dv'$ and integrating we get that
\begin{align*}
    \frac{dv}{1 - v^2/v^2} &= -\frac{v_{ex}}{M}dM \\
    \int_0^v \frac{dv}{1 - v^2/v^2} &= -v_{ex}\int_{M_i}^M\frac{dM}{M} \\
    c\arctan(\frac{v}{c}) &= -v_{ex} [\log(M) - \log(M_i)] \\
    \frac{c}{v_{ex}}\arctan(\frac{v}{c}) &= \log(\frac{M_i}{M}) \\
    \frac{c}{2v_{ex}}\log(\frac{1 + v/c}{1 - v/c}) &= \log(\frac{M_i}{M})
\end{align*}
Therefore
\begin{align*}
    \log(\frac{M_i}{M}) &= \log(\bigg(\frac{c + v}{c - v}\bigg)^{\frac{c}{2v_{ex}}}) \\
    \frac{M_i}{M} &= \bigg(\frac{c + v}{c - v}\bigg)^{\frac{c}{2v_{ex}}}
\end{align*}
\item [(b)] Let $v_{ex} = 3\sci{5}~cm/s$, $v = 0.1 c$ and $M = 1~ton$ then
the mass of fuel needed is 
\begin{align*}
    M_i - M &= M\bigg(\frac{c + v}{c - v}\bigg)^{\frac{c}{2v_{ex}}} - M \\
    &=\bigg(\frac{1.1}{0.9}\bigg)^{\frac{3\sci{10}}{6\sci{5}}} - 1 \\
    &= 3.12 \sci{4354}~ton
\end{align*}
\item [(c)] Let $v_{ex} = c$, $v = 0.1 c$ and $M = 1~ton$ then
the mass of fuel needed is 
\begin{align*}
    M_i - M &= M\bigg(\frac{c + v}{c - v}\bigg)^{\frac{c}{2v_{ex}}} - M \\
    &=\bigg(\frac{1.1}{0.9}\bigg)^{1/2} - 1 \\
    &= 0.105~ton
\end{align*}

\end{itemize}
\end{proof}

\cleardoublepage
\begin{proof}[\proofname\ 22]
Let $\Phi(t,\bm{x})$ be a scalar field. To show that 
\begin{align*}
    f(x^\mu) = \int \Phi(x^\mu)~dVdt
\end{align*}
is a scalar, we must show that $f$ is invariant under Lorentz transformations.
If $f(x^\mu)$ is a scalar then must be that $f(x^\mu) = f(a_\alpha^\mu x'^\alpha)$. \\
Let us compute $f(a_\alpha^\mu x'^\alpha)$ as follows
\begin{align*}
    f(a_\alpha^\mu x'^\alpha)
    &= \int \Phi(a_\alpha^\mu x'^\alpha)~dV'dt' \\
    &= \int \Phi(x^\mu)~\sqrt{1 - V^2}dV~\frac{dt}{\sqrt{1 - V^2}} \\
    &= \int \Phi(x^\mu)~dVdt \\
    &= f(x^\mu)
\end{align*}
We used that $\Phi$ is a scalar field and hence
$\Phi(a_\alpha^\mu x'^\alpha) = \Phi(x^\mu)$, also, we used the volume
contraction and time dilation equations.
\end{proof}

\cleardoublepage
\begin{proof}[\proofname\ 23]
Let the current density $j^\mu$ differ from zero only in a volume of
finite extent in the $x, y, z$ directions.\\
Let us integrate $\partial_\mu j^\mu = 0$ over the volume of a four-dimensional
world tube that encloses the region of interest and extends from $x^0 =$
constant to $x'^0 =$ constant.
\begin{align*}
    \int_\Omega \partial_\mu j^\mu ~ d\Omega &= 0
\end{align*}
We note that we are integrating over a four-dimensional volume $\Omega$,
then applying Gauss' theorem we get that
\begin{align*}
    \int_V j^\mu n_\mu ~ dV &= 0
\end{align*}
Where $V$ is the surface of our four-dimensional volume. Given that $j^\mu$ is
only non-zero in a finite volume which is enclosed by the integral, then
the "flux" is only non-zero in the direction of the axis $x^0$, so the integral
becomes
\begin{align*}
    \int_V j'^0 ~ dV - \int_V j^0 ~ dV &= 0
\end{align*}
Therefore 
\begin{align*}
    Q' = \int_V j'^0 ~ dV = \int_V j^0 ~ dV = Q
\end{align*}
This implies that the total charge is a scalar under Lorentz transformations.
\end{proof}

\cleardoublepage
\begin{proof}[\proofname\ 24]
In the previous problem we found that
\begin{align*}
    Q = \int_V J^0 dV
\end{align*}
Since $J^\mu$ is different from zero only in a volume of a finite extent then
the exchange of charge with the space outside the volume is zero, then must be
that the total charge $Q$ held inside the volume doesn't change with time, 
therefore
\begin{align*}
    \dv{Q}{t} = \dv{t}(\int_V J^0 dV) = 0
\end{align*}
\end{proof}

\cleardoublepage
\begin{proof}[\proofname\ 25]
In Exercise 17 we showed that an electrically neutral wire carring a current
in the $xyz$ reference frame is not electrically neutral in the reference 
frame $x'y'z'$ and the charge density (in this case a linear charge density)
is given by
\begin{align*}
    \lambda' = -\frac{Vj_x/c^2}{\sqrt{1 - V^2/c^2}}
\end{align*}
So if $j_x = 3\sci{11}~esu/s$ and $V = 0.8c$ we get that the charge per unit of
length is
\begin{align*}
    \lambda' = -\frac{0.8 \cdot 3\sci{11}/3\sci{10}}{\sqrt{1 - 0.8^2}}
    = -13.33~esu/cm
\end{align*}
\end{proof}

\cleardoublepage
\begin{proof}[\proofname\ 26]
We know that an electrically neutral wire carring a current
in the $xyz$ reference frame is not electrically neutral in the reference 
frame $x'y'z'$ and the charge density (in this case a linear charge density)
is given by
\begin{align*}
    \lambda' = -\frac{VI_x/c^2}{\sqrt{1 - V^2/c^2}}
\end{align*}
Where $I_x$ is the $x$ component of the current $I$, in this case $I_x$ is
related to $\theta$ by
\begin{align*}
    I_x = I\cos\theta
\end{align*}
Therefore the charge density observed in the reference frame $x'y'z'$ is
\begin{align*}
    \lambda' = -\frac{VI\cos\theta/c^2}{\sqrt{1 - V^2/c^2}}
\end{align*}
\end{proof}

\cleardoublepage
\begin{proof}[\proofname\ 27]
Let us consider two large, horizontal, parallel conducting plates carring
currents in opposite directions, the positive $x$ direction and the negative
$x$ direction. The current is uniformly distributed over each plate, with $k$
esu of charge per second flowing across each centimeter measured along the
plate perpendicularly to the current.\\
An observer in a reference frame moving at speed $V$ in the positive $x$
direction will see a charge density and a current density of
\begin{align*}
    \sigma'_{+x} = -\frac{Vk}{c^2\sqrt{1 - V^2/c^2}}
    \qquad
    j_{+x}' = \frac{k}{\sqrt{1 - V^2/c^2}}
\end{align*}
for the plate with current in the positive $x$ direction.\\
For the plate with current in the negative $x$ direction, the observer will
experience a charge density and a current density of
\begin{align*}
    \sigma'_{-x} = \frac{Vk}{c^2\sqrt{1 - V^2/c^2}}
    \qquad
    j_{-x}' = -\frac{k}{\sqrt{1 - V^2/c^2}}
\end{align*}
We know that the electric field due to two charged plates is zero outside the
plates (they cancel out) and in between them is 
\begin{align*}
    E_z = 4\pi \sigma
\end{align*}
Therefore
\begin{align*}
    E_z = \frac{4\pi Vk}{c^2\sqrt{1 - V^2/c^2}}
\end{align*}
\end{proof}
\end{document}

